% Generated draft: requires mathematical and source review.
\documentclass[11pt]{article}
\usepackage[margin=25mm]{geometry}
\usepackage{fontspec}
\setmainfont{Latin Modern Roman}
\usepackage{amsmath,amssymb,mathtools,unicode-math}
\setmathfont{Latin Modern Math}
\usepackage{enumitem,needspace,xurl,hyperref}
\setlength{\parindent}{0pt}
\setlength{\parskip}{5pt}
\setlength{\emergencystretch}{3em}
\newcommand{\N}{\mathbb{N}}
\newcommand{\Z}{\mathbb{Z}}
\newcommand{\Q}{\mathbb{Q}}
\newcommand{\R}{\mathbb{R}}
\newcommand{\C}{\mathbb{C}}
\newcommand{\F}{\mathbb{F}}
\newcommand{\A}{\mathbb{A}}
\newcommand{\PP}{\mathbb{P}}
\newcommand{\E}{\mathbb{E}}
\newcommand{\eps}{\varepsilon}
\newcommand{\dd}{\,\mathrm{d}}
\begin{document}
\section*{UnsolvedMath ID: 701686}
\textbf{Source problem code:} catalogue entry
\subsection{Definitions and mathematical statement}
For a finite group \(K\), define its set of conjugacy class sizes by
\[
\operatorname{cs}(K)=\{\lvert x^K\rvert:x\in K\}
=\{[K:C_K(x)]:x\in K\},
\qquad C_K(x)=\{k\in K:kx=xk\}.
\]
This is a set without multiplicities and includes \(1\).

Let \(L\) be a non-abelian finite simple group. Write \(M(L)=H_2(L,\Z)\) for its Schur multiplier, and let \(H(L)\) be its universal perfect central extension. Thus \(H(L)=[H(L),H(L)]\), and there is an exact sequence
\[
1\longrightarrow M(L)\longrightarrow H(L)\xrightarrow{\pi}L\longrightarrow 1,
\qquad \ker\pi\leq Z(H(L)),
\]
with the following universal property: for every central extension \(q:E\twoheadrightarrow L\), there is a unique homomorphism \(f:H(L)\to E\) such that \(q\circ f=\pi\). The source notation \(M(L).L\) denotes this extension and does not assert that it splits.

\textbf{Question.} For every non-abelian finite simple group \(L\) and every finite group \(G\), does
\[
\operatorname{cs}(G)=\operatorname{cs}(H(L))
\quad\Longrightarrow\quad
\exists\,A\text{ finite abelian such that }G\cong H(L)\times A
\]
hold?

The supplied source reports that this assertion has been proved when \(L=\operatorname{Alt}_5\), the alternating group on five letters. That special-case result is a supplied literature claim; it has not been independently verified here, and its supplied citation lacks authors and an article title.
\subsection{Short English statement}
Does the set of conjugacy class sizes determine the universal perfect central extension of every non-abelian finite simple group up to a direct product with a finite abelian group? The supplied record labels the general question open and reports a proof for \(L=\operatorname{Alt}_5\).
\subsection{Sources}
\begin{itemize}
\item Local UnsolvedMath registry, supplied complete source record: permanent ID \(701686\), exact source problem code \texttt{catalogue entry}, title \emph{Kourovka Notebook Problem 20.28}. The record labels the problem \emph{open}; this is supplied registry metadata, not an independently verified status assessment. Upstream revision: \texttt{d371ee917ac9647fbb06cd6fc529deec6f3dda82}; record updated \texttt{2026-10-16T14:00:00Z}.
\item I. B. Gorshkov, proposer, \emph{The Kourovka Notebook (Unsolved Problems in Group Theory)}, Problem 20.28, 20th issue (2022). The supplied background also identifies the 21st issue edition (2026), \texttt{arXiv:1401.0300v48}. The supplied collection description names E. I. Khukhro and V. D. Mazurov as editors.
\item Supplied reference for the case \(L=\operatorname{Alt}_5\): \emph{J. Algebra Appl.}, volume \textbf{21}, no. 11 (2022), Article ID 2250226, 8 pages. Authors, article title, and a link are not supplied.
\end{itemize}
\end{document}
